\chapter{CONCLUSIONES}
\section{Conclusiones del trabajo}

El objetivo principal de este Trabajo de Fin de Grado era estudiar la aplicación de las \textit{Tensor Networks} en modelos de \textit{deep learning} para clasificación de imágenes. A partir del trabajo realizado, se puede concluir que este objetivo se ha cumplido de forma satisfactoria. El proyecto no se ha limitado a una revisión teórica de las redes tensoriales, sino que ha incluido el diseño, implementación, entrenamiento y evaluación de un modelo tensorial integrado en un \textit{pipeline} completo de clasificación.

En la parte teórica, el trabajo ha permitido establecer una base sólida sobre los conceptos necesarios para comprender este tipo de modelos, entre ellos los tensores, las contracciones, las descomposiciones tensoriales, la notación diagramática y topologías como MPS, TTN, PEPS o MERA. Esta fundamentación ha resultado clave para justificar las decisiones tomadas más adelante en la implementación, en especial la elección de una arquitectura jerárquica tipo TTN como núcleo del modelo desarrollado.

En el terreno computacional, se ha conseguido construir un sistema entrenable y reproducible basado en PyTorch y TensorKrowch. El proyecto abarca la carga y preparación de datos, el preprocesamiento de imágenes, la generación de recortes a partir de anotaciones YOLO, la definición de experimentos mediante archivos de configuración, el entrenamiento con GPU, el registro de métricas y la generación de diagnósticos. Esta organización ha permitido comparar configuraciones de manera ordenada y seguir la evolución del modelo a lo largo del desarrollo.

Una conclusión relevante es que las \textit{Tensor Networks}, concretamente para clasificación de imágenes, necesitan tener un buen tratamiento de datos e información, para que cuando entren en el modelo, le resulte más fácil diferenciar entre clases y mejorar las predicciones. Los primeros experimentos mostraron que una arquitectura tensorial aplicada directamente sobre representaciones simples tenía dificultades para generalizar. La incorporación de \texttt{AIMPatchEmbedding} permitió conservar mejor la información espacial de la imagen antes de entregarla a la TTN, y se convirtió en uno de los cambios clave para mejorar el comportamiento del entrenamiento.

El cambio al \textit{dataset} ChemEq25 también fue decisivo. A diferencia de Flowers102, que presentaba un número elevado de clases visualmente muy parecidas, ChemEq25 permitió abordar un problema más alineado con los objetivos del proyecto, la clasificación de material de laboratorio a partir de recortes individuales.

Los resultados experimentales muestran una mejora progresiva y coherente. El modelo pasó de unos resultados iniciales limitados a alcanzar un rendimiento elevado tras introducir mejoras en la representación de entrada, la optimización, la regularización y el \textit{data augmentation}. El mejor modelo, basado en la combinación de \texttt{AIMPatchEmbedding} y \texttt{HybridCustomTTN}, obtuvo un 96,05\% de \textit{accuracy}, un 95,94\% de F1 macro y un 99,43\% de \textit{top-5 accuracy} sobre ChemEq25. Estos valores indican que la arquitectura propuesta es capaz de aprender patrones visuales relevantes.

También se ha comprobado que aumentar la capacidad del modelo no garantiza por sí solo una mejor generalización. En los experimentos finales, los cambios que más contribuyeron al rendimiento fueron la preservación de la relación de aspecto, el \textit{label smoothing}, la regularización y, sobre todo, una estrategia suave de \textit{data augmentation}. Esto sugiere que, en modelos tensoriales aplicados a imágenes, la preparación de los datos y el control del entrenamiento pueden llegar a ser tan importantes como la propia topología de la red.

Entre las limitaciones, conviene señalar que no se ha realizado una comparación exhaustiva frente a arquitecturas convencionales como las CNN o los Transformers, ya que el objetivo principal era estudiar la viabilidad de las redes tensoriales, una línea que podría desarrollarse en trabajos posteriores. Asimismo, se ha observado cierta sensibilidad a la \textit{seed} y a la configuración de entrenamiento, lo que indica que estos modelos requieren una validación experimental cuidadosa.

En conjunto, el trabajo demuestra que las \textit{Tensor Networks} pueden integrarse de manera efectiva en modelos de \textit{deep learning} para clasificación de imágenes. Su aplicación práctica exige combinar la estructura tensorial con técnicas habituales del \textit{deep learning}, una representación visual adecuada y un seguimiento riguroso de los experimentos. La principal aportación del proyecto es, por tanto, mostrar un recorrido completo que va desde los fundamentos físicos y matemáticos de las redes tensoriales hasta un modelo funcional, entrenable y evaluable sobre datos reales.


\section{Conclusiones personales}

El desarrollo de este Trabajo de Fin de Grado ha supuesto una experiencia especialmente enriquecedora, tanto desde el punto de vista académico como personal. En primer lugar, el tema abordado pertenece a un ámbito muy especializado, las \textit{tensor networks}, cuyo funcionamiento requiere una comprensión profunda tanto de conceptos físicos como de herramientas computacionales avanzadas. Haber tenido la oportunidad de trabajar en este contexto, y de aprender bajo la supervisión de investigadores con experiencia en el área, en particular Alejandro Mata Ali y todo el departamento de computación cuántica del Instituto Tecnológico de Castilla y León, ha sido un privilegio y una parte fundamental del proceso formativo asociado a este trabajo.

A lo largo de estos meses, este TFG no solo me ha permitido profundizar en la física de las \textit{tensor networks}, sino también adquirir una gran cantidad de competencias técnicas que han resultado esenciales para el desarrollo del proyecto. Entre ellas, destaca el aprendizaje de Python a un nivel mucho más avanzado, así como la experiencia de programar desde cero un proyecto de cierta envergadura. Esto ha requerido no solo implementar código funcional, sino también aprender a organizarlo correctamente, aplicar buenas prácticas de programación y estructurar el trabajo de forma clara y mantenible.

Además, el proyecto ha supuesto un reto importante en el ámbito del \textit{machine learning}. Ha sido necesario familiarizarse con PyTorch, estudiar las bases matemáticas del \textit{deep learning} y comprender el flujo general de trabajo asociado a este tipo de modelos, desde la preparación de los datos hasta el entrenamiento y análisis de resultados. Todo ello ha ampliado considerablemente mi formación computacional y me ha permitido conectar herramientas modernas de inteligencia artificial con problemas propios de la física teórica y computacional.

Otro aspecto fundamental del desarrollo de este trabajo ha sido el aprendizaje de herramientas habituales en el entorno investigador. Durante el proyecto he aprendido a utilizar Git y GitHub para el control de versiones y la gestión del código, así como comandos básicos de Linux y Bash necesarios para trabajar de manera eficiente en los servidores de la universidad (Lorca). También he tenido la oportunidad de utilizar recursos de computación de alto rendimiento, que han sido clave para llevar a cabo los cálculos y entrenamientos necesarios. Esta experiencia me ha permitido comprender mejor cómo se organiza y ejecuta un proyecto computacional real en un entorno de investigación.

Más allá de los conocimientos técnicos, este TFG me ha acercado al funcionamiento cotidiano de la investigación científica. Haber podido participar en reuniones, observar la forma de trabajo del grupo y ver de primera mano cómo se plantean, discuten y desarrollan problemas de investigación ha sido una parte especialmente valiosa de la experiencia. En este sentido, el trabajo no solo ha consistido en resolver un problema concreto, sino también en aprender cómo se investiga, cómo se colabora y cómo se construye conocimiento en un contexto académico real.

En conjunto, este TFG ha representado un desafío constante. Por un lado, debido a la dificultad conceptual asociada a la física de las \textit{tensor networks}; por otro, por la exigencia técnica derivada de la programación, el uso de herramientas de aprendizaje automático y el trabajo en entornos computacionales avanzados. Sin embargo, precisamente esta combinación de retos ha hecho que el proceso sea especialmente formativo. En apenas unos meses he adquirido conocimientos y habilidades que al inicio del proyecto no imaginaba poder desarrollar, y considero que esta experiencia ha supuesto un punto de inflexión importante en mi formación como físico.

En definitiva, ha sido un proyecto muy ambicioso que ha puesto a prueba todos mis conocimientos sobre física aprendidos en la carrera y ha ampliado considerablemente esos conocimientos, mas especializados en la física computacional. También me ha permitido crecer a nivel científico, ya que he aprendido muchas herramientas que me servirán en mi carrera profesional.